\section{Introduction}
Variational inference approximates a posterior by optimizing a tractable density.
Finite mixtures make the approximation more expressive by representing several
modes with separate components. Variational boosting adds these components
incrementally \citep{miller2017boosting}, while automatic differentiation
variational inference with mixtures jointly optimizes their parameters and weights
\citep{morningstar2021mixtures}.

Sampling a mixture involves drawing a categorical index and then a continuous
value from the chosen component. Training requires expectation gradients that
account for parameter changes in both stages. Score-function estimators use
derivatives of the log density \citep{ranganath2014bbvi}, while pathwise estimators
differentiate the sampled value, as in variational autoencoders
\citep{kingma2014vae,rezende2014backprop}. Each approach has its own variance
and computational requirements \citep{mohamed2020gradients}.

Several methods adapt these ideas to mixtures. Stratification samples every
component and combines the resulting gradients using the mixture weights
\citep{roeder2017stl,morningstar2021mixtures}. Implicit reparameterization obtains
sample derivatives from distribution functions
\citep{graves2016mixture,figurnov2018implicit}. Transport-based methods describe
how probability mass moves as parameters change
\citep{jankowiak2018transport,jankowiak2019multivariate}. This perspective also
leads to a post-stratified estimator, which averages component derivatives at a
single mixture sample. \citet{burroni2026mixture} bring these approaches into a
common formulation and analyze their computational and variance properties.

We will implement four algorithms as the ReMix library: a score-function
baseline, stratification, implicit reparameterization, and post-stratification.
The library will provide a shared PyTorch interface, starting with mixtures of
diagonal Gaussian components, and support the Sticking the Landing (STL)
modification for evidence lower bound gradients \citep{roeder2017stl}.
This report specifies the estimators and their implementation requirements as a
basis for comparing gradient variance, optimization progress, and runtime.
